\section{有限样本假设检验}

估计给出一个参数值或区间；检验则预先划定哪些数据足以反对零假设。若在看见数据后才改变拒绝规则，原先设定的错误率一般不再成立。

\begin{definition}[检验、大小与功效]
检验函数 \(\varphi(x)\in[0,1]\) 表示观测 \(x\) 后拒绝零假设的概率；取 \(0\) 或 \(1\) 时是非随机化检验。对 \(H_0:\theta\in\Theta_0\)，其大小为
\(\sup_{\theta\in\Theta_0}\mathbb E_\theta\varphi(X)\)；
在备择 \(\theta\) 下的功效是 \(\mathbb E_\theta\varphi(X)\)。
所谓水平 \(\alpha\)，指大小不超过 \(\alpha\)。
\end{definition}

\begin{theorem}[Neyman--Pearson 引理，有限样本]
设简单零假设与备择的密度 \(p_0,p_1\) 具有共同支配测度。选 \(c\ge0\)，令
\(\varphi^*(x)=1\) 当 \(p_1(x)>cp_0(x)\)，
\(\varphi^*(x)=0\) 当 \(p_1(x)<cp_0(x)\)，
等号处可随机化，使 \(\mathbb E_0\varphi^*=\alpha\)。
则对任何满足 \(\mathbb E_0\varphi\le\alpha\) 的检验 \(\varphi\)，有
\(\mathbb E_1\varphi\le\mathbb E_1\varphi^*\)。
\end{theorem}
\begin{proof}
在 \(p_1-cp_0>0\) 处，\(\varphi^*=1\ge\varphi\)；在
\(p_1-cp_0<0\) 处，\(\varphi^*=0\le\varphi\)。
因此逐点有
\((\varphi^*-\varphi)(p_1-cp_0)\ge0\)。
积分后得
\[
\mathbb E_1(\varphi^*-\varphi)
\ge c\mathbb E_0(\varphi^*-\varphi)
\ge0,
\]
末步使用检验大小约束。等号集合的随机化只为精确达到指定大小。
\end{proof}

例如 \(X\sim\operatorname{Bin}(n,p)\)，检验
\(H_0:p=p_0\) 对 \(H_1:p=p_1>p_0\)。似然比为
\((p_1/p_0)^X[(1-p_1)/(1-p_0)]^{n-X}\)，随 \(X\) 严格递增；
所以拒绝域形如 \(X>k\)，必要时在 \(X=k\) 随机化。阈值由
\(P_{p_0}(X>k)\le\alpha\) 定，功效则在 \(P_{p_1}\) 下计算。一个 \(p\) 值是零假设下预先选定的“至少如此极端”的尾概率，并不是零假设为真的概率。

这个例子还允许把复合假设的一种特殊情形证明到底。考虑
\(H_0:p\le p_0\) 对 \(H_1:p>p_0\)，用同一个单调拒绝函数
\(\varphi(X)\)，令其在 \(p_0\) 下大小恰为 \(\alpha\)。取独立均匀变量 \(U_i\sim\operatorname{Uniform}(0,1)\)，构造
\(X(p)=\sum_{i=1}^n\mathbf1_{\{U_i\le p\}}\)。对每条样本路径，\(p\) 增大时 \(X(p)\) 不减；所以单调函数
\(\varphi(X(p))\) 的期望也不减。于是
\(\sup_{p\le p_0}\mathbb E_p\varphi=\mathbb E_{p_0}\varphi=\alpha\)，即控制复合零假设的大小。另一方面，对每个固定 \(p_1>p_0\)，比值
\(p_{p_1}(x)/p_{p_0}(x)\) 都随 \(x\) 增加，Neyman--Pearson 引理说同一阈值检验在所有大小为 \(\alpha\) 的检验中对该 \(p_1\) 功效最大。由于阈值不依赖所选 \(p_1\)，它对整个单侧备择\emph{一致最强}。这里的结论依赖单调似然比；双侧备择一般没有同一个阈值可以同时照顾两边。

最小的数值例子能看清随机化的必要性。只观察一次，取
\(p_0=0.3\)、\(p_1=0.8\)、\(\alpha=0.1\)。若见成功就必拒绝，零假设下拒绝概率为 \(0.3\)，超过目标水平；若永不拒绝，功效又为零。在 \(X=1\) 时以概率 \(1/3\) 拒绝，在 \(X=0\) 时不拒绝，大小为
\(0.3/3=0.1\)，在备择下功效为 \(0.8/3\)。随机化来自离散样本空间无法用确定性阈值精确切出任意 \(\alpha\)，而不是对观察结果的含糊处理。

简单假设的最优性不能直接推广到任意复合假设。若 \(H_0\) 或 \(H_1\) 包含多个参数值，需额外证明某检验对所有相关参数同时最优，或者接受一种较弱的最优标准。

\begin{exercise}
只观察一次 Bernoulli\((p)\) 试验，检验 \(p=p_0\) 对 \(p=p_1>p_0\)。当 \(0<\alpha<p_0\) 时，求在 \(X=1\) 上的随机化概率与检验功效。
\end{exercise}
